\subsection{全ての pair が非零である四ray core の正符号グラフ}
\label{fg24:section}

三高と四高に共通する短い整数論的な正値性判定を与える．
以下の結果は，四高 core に三高の完全マッチング定理をそのまま適用するものではない．
四高の場合には道型 $P_4$ を必要候補として残す．

\begin{lemma}[正符号の孤立頂点は存在しない]
\label{fg24:no-isolated}
四つの正整数 $q_i$ と整数 $p_{ij}$ が
\[
 \sum_{j\ne i}p_{ij}=-q_i,\qquad
 p_{ij}p_{ji}=q_iq_j-1>0
\]
を満たすとする．さらに正の basin weights が存在するなら，
$p_{ij}>0$ によって定まる無向グラフ $G^+$ は孤立頂点を持たない．
\end{lemma}

\begin{proof}
頂点 $h$ が孤立していると仮定し，残る三頂点について
$c_i=-p_{hi}>0$ と置く．行和から $\sum_{i\ne h}c_i=q_h$ である．
異なる $i,j\ne h$ に対し，整数
\[
 g_{ij}=q_i c_j-c_i p_{ij}
\]
を定める．行和と pair products を使うと
\begin{align*}
 \sum_{j\ne i,h}g_{ij}
 &=q_i(q_h-c_i)-c_i(-q_i-p_{ih})\\
 &=q_iq_h+c_i p_{ih}=1.
\end{align*}
左辺は二つの整数の和なので，少なくとも一つは非正である．
そのような $j$ を各 $i$ について一つ選ぶと，$p_{ij}>0$ かつ
\[
 \frac{p_{ij}}{q_i}\ge\frac{c_j}{c_i}
\]
となる．三頂点上で各頂点から一つずつ辺を選んだ有向グラフは，
有向閉路を持つ．閉路に沿って上式を掛けると，その gain は一以上になる．
これは正の weights による狭義の閉路不等式
\eqref{kl10:cycle-test} に反する．
\end{proof}

\begin{theorem}[非零 pair の四ray正値性判定]
\label{fg24:criterion}
前補題の行和・整数性・狭義正の pair products を仮定する．
次の三条件は同値である．
\begin{enumerate}
 \item 正の basin weights が存在する．
 \item $G^+$ は孤立頂点を持たない．
 \item $G^+$ は完全マッチング $2K_2$ または道 $P_4$ である．
\end{enumerate}
従って三高・四高 core の正値性は，この二つの符号型で判定できる．
四高の $P_4$ 型を持つ整数模型の存在は，この定理の主張に含めない．
\end{theorem}

\begin{proof}
第一条件から第二条件は前補題である．行和が負なので，各頂点には少なくとも
一本の負の辺がある．従って $G^+$ の次数は二以下である．
第二条件の下では全ての次数が一または二となり，四頂点上の候補は
$2K_2,P_4,C_4$ だけである．
$C_4$ の場合，補グラフはマッチングである．その負の辺 $ij$ の両端は
他の二つの entries が正なので，行和から
\[
 -p_{ij}>q_i,\qquad -p_{ji}>q_j
\]
となる．従って $p_{ij}p_{ji}>q_iq_j$ となり，pair product と矛盾する．
これで第三条件を得る．

第三条件のグラフはいずれも森である．正の simple directed cycles は
長さ二だけであり，その gain は
$(q_iq_j-1)/(q_iq_j)<1$ である．従って閉路判定から第一条件が従う．
正の有理 weights を得るには，各木の根の weight を一とし，辺に沿って
\[
 \frac{p_{ij}}{q_i}<\frac{w_j}{w_i}<\frac{q_j}{p_{ji}}
\]
の開区間内の有理数を順次選んでもよい．木なので整合性の衝突はない．
\end{proof}

\begin{remark}[整数性と適用範囲]
この証明の要点は，二つの\emph{整数}の和が一なら一方は非正になることである．
実数 entries に同じ議論を適用してはならない．また二本以上の低 ray があると
零の pair product が生じ得るため，その場合は仮定の外である．
正の局所模型の判定であり，整数 clutching，node 配分，対向 $D$ 相殺や
大域 atlas の存在を与えるものではない．
\end{remark}

\paragraph{三高の既存証明への応用．}
低 ray を含む三角形と孤立頂点からなる符号型は，補題~\ref{fg24:no-isolated}
で直ちに排除される．三高の場合の二つの道型の整数障害と合わせると，
定理~\ref{mh18:matching} の完全マッチング結論を得る．
四高では，三高の道型排除に使う低分母一の仮定が失われるため，
$P_4$ の可能性を別問題として保持する．
